\documentclass[11pt,a4paper]{article}
\usepackage[margin=19mm,top=17mm,bottom=17mm]{geometry}
\usepackage{amsmath,amssymb,amsfonts}
\usepackage{mathptmx}
\usepackage{enumitem}
\usepackage{fancyhdr}
\usepackage{array,tabularx,booktabs}
\usepackage[hidelinks]{hyperref}
\setlength{\parindent}{0pt}\setlength{\parskip}{4pt}
\setlist[enumerate,1]{leftmargin=1.1em,itemsep=3pt,parsep=2pt}
\setlist[itemize]{leftmargin=1.4em,itemsep=1pt}
\pagestyle{fancy}\fancyhf{}
\fancyhead[C]{\footnotesize \textbf{Question Paper}}
\fancyfoot[C]{\footnotesize Page \thepage}
\renewcommand{\headrulewidth}{0.4pt}
\begin{document}
\vspace{2mm}
\noindent \textbf{Marks : 94} \hfill \textbf{Time : 2 : 30 hours}
\noindent\rule{\linewidth}{1.1pt}\vspace{-1mm}\noindent\rule{\linewidth}{0.5pt}
{\centering \Large \textbf{INDEFINITE INTEGRALS TEST PAPER}\par}
{\centering JEE $|$ Mathematics\par}
\noindent\rule{\linewidth}{1.1pt}\vspace{-1mm}\noindent\rule{\linewidth}{0.5pt}
\vspace{1mm}
\noindent \textbf{Duration:} Not specified \quad \textbf{Total Marks:} 94 \quad \textbf{Questions:} 25\par
\vspace{2mm}

\begin{tabularx}{\linewidth}{|X|X|}\hline
Candidate Name: \rule{6cm}{0.4pt} & Roll Number: \rule{3.5cm}{0.4pt} \\\hline
\end{tabularx}

\textbf{Instructions}
\begin{itemize}[topsep=2pt]
\item Read every question carefully before selecting an answer.
\item For single-correct questions, mark one option only.
\item Use the answer sheet or space provided by the invigilator.
\end{itemize}
\begin{enumerate}[resume]
\item Suppose \(\displaystyle \int x^{4}e^{-2x}\,dx=e^{-2x}f(x)+C\), where \(C\) is a constant and \(f(x)\) is a polynomial. Then the value of \(f(1)\) is \hfill [4 marks]
{\renewcommand{\arraystretch}{1.8}
\begin{tabularx}{\linewidth}{X X X X}
(A) \quad \(-\dfrac{21}{4}\) & (B) \quad \(-\dfrac{15}{4}\) & (C) \quad \(\dfrac{21}{4}\) & (D) \quad \(-\dfrac{9}{4}\) \\[10pt]
\end{tabularx}}
\item Let $g(x)$ be a polynomial such that
\[
\int x^3e^{-2x}\,dx=g(x)e^{-2x}+C,
\]
where $C$ is an arbitrary constant. Then the value of $g(-1)$ is \hfill [4 marks]
{\renewcommand{\arraystretch}{1.8}
\begin{tabularx}{\linewidth}{X X X X}
(A) \quad $-\dfrac{1}{8}$ & (B) \quad $\dfrac{1}{8}$ & (C) \quad $\dfrac{3}{8}$ & (D) \quad $-\dfrac{3}{8}$ \\[10pt]
\end{tabularx}}
\item Let
\[
\int x^{4}e^{-2x}\,dx=e^{-2x}g(x)+C,
\]
where \(C\) is an arbitrary constant and \(g(x)\) is a polynomial. Then the value of \(g(1)\) is \hfill [4 marks]
{\renewcommand{\arraystretch}{1.8}
\begin{tabularx}{\linewidth}{X X X X}
(A) \quad \(-\dfrac{15}{4}\) & (B) \quad \(-\dfrac{21}{4}\) & (C) \quad \(-\dfrac{23}{4}\) & (D) \quad \(-\dfrac{25}{4}\) \\[10pt]
\end{tabularx}}
\item Let \(g(x)\) be a polynomial of degree at most \(4\) such that
\[
\int x^4e^{-2x}\,dx=e^{-2x}g(x)+C,
\]
where \(C\) is a constant. Then \(g(0)\) is equal to: \hfill [4 marks]
{\renewcommand{\arraystretch}{1.8}
\begin{tabularx}{\linewidth}{X X X X}
(A) \quad \(-\dfrac{3}{4}\) & (B) \quad \(\dfrac{3}{4}\) & (C) \quad \(-\dfrac{3}{2}\) & (D) \quad \(\dfrac{1}{2}\) \\[10pt]
\end{tabularx}}
\item Let \(G(x)\) be a polynomial such that
\[
\int x^4 e^{-2x}\,dx=e^{-2x}G(x)+C,
\]
where \(C\) is an arbitrary constant. The value of \(G(1)\) is \hfill [4 marks]
{\renewcommand{\arraystretch}{1.8}
\begin{tabularx}{\linewidth}{X X X X}
(A) \quad \(-\dfrac{21}{4}\) & (B) \quad \(-\dfrac{19}{4}\) & (C) \quad \(\dfrac{21}{4}\) & (D) \quad \(-\dfrac{25}{4}\) \\[10pt]
\end{tabularx}}
\item If \(\displaystyle \int x^{3}e^{2x}\,dx=e^{2x}f(x)+C\), where \(C\) is a constant of integration and \(f(x)\) is a polynomial, then the value of \(f(-1)\) is \hfill [4 marks]
{\renewcommand{\arraystretch}{1.8}
\begin{tabularx}{\linewidth}{X X X X}
(A) \quad \(-\dfrac{29}{16}\) & (B) \quad \(-\dfrac{21}{16}\) & (C) \quad \(\dfrac{29}{16}\) & (D) \quad \(\dfrac{21}{16}\) \\[10pt]
\end{tabularx}}
\item Let \(G(x)\) be a polynomial such that
\[
\int x^{4}e^{-2x}\,dx=e^{-2x}G(x)+C,
\]
where \(C\) is an arbitrary constant. The value of \(G(-1)\) is \hfill [3 marks]
{\renewcommand{\arraystretch}{1.8}
\begin{tabularx}{\linewidth}{X X X X}
(A) \quad \(-\dfrac{1}{4}\) & (B) \quad \(\dfrac{1}{4}\) & (C) \quad \(-\dfrac{3}{4}\) & (D) \quad \(\dfrac{3}{4}\) \\[10pt]
\end{tabularx}}
\item Let $g(x)$ be a polynomial of degree at most $3$. If
\[
\int x^3e^{-2x}\,dx=e^{-2x}g(x)+C,
\]
where $C$ is a constant of integration, then the value of $g(1)$ is \hfill [4 marks]
{\renewcommand{\arraystretch}{1.8}
\begin{tabularx}{\linewidth}{X X X X}
(A) \quad $-\dfrac{19}{8}$ & (B) \quad $-\dfrac{13}{8}$ & (C) \quad $\dfrac{19}{8}$ & (D) \quad $-\dfrac{11}{8}$ \\[10pt]
\end{tabularx}}
\item Let \(g(x)\) be a polynomial such that
\[
\int x^{4}e^{-2x}\,dx=e^{-2x}g(x)+C,
\]
where \(C\) is an arbitrary constant. Then the value of \(g(1)\) is \hfill [4 marks]
{\renewcommand{\arraystretch}{1.8}
\begin{tabularx}{\linewidth}{X X X X}
(A) \quad \(-\dfrac{17}{4}\) & (B) \quad \(-\dfrac{21}{4}\) & (C) \quad \(-\dfrac{23}{4}\) & (D) \quad \(-\dfrac{25}{4}\) \\[10pt]
\end{tabularx}}
\item Let \(f(x)\) be a polynomial such that
\[
\int x^4e^{-2x}\,dx=e^{-2x}f(x)+C,
\]
where \(C\) is a constant of integration. Then the value of \(f(1)\) is \hfill [4 marks]
{\renewcommand{\arraystretch}{1.8}
\begin{tabularx}{\linewidth}{X X X X}
(A) \quad \(-\dfrac{15}{4}\) & (B) \quad \(-\dfrac{21}{4}\) & (C) \quad \(\dfrac{21}{4}\) & (D) \quad \(-\dfrac{25}{4}\) \\[10pt]
\end{tabularx}}
\item Let \(\int x^{4}e^{-2x}\,dx=e^{-2x}g(x)+C\), where \(g(x)\) is a polynomial and \(C\) is an arbitrary constant. Then the value of \(g(1)\) is \hfill [1 marks]
{\renewcommand{\arraystretch}{1.8}
\begin{tabularx}{\linewidth}{X X X X}
(A) \quad \(-\frac{21}{4}\) & (B) \quad \(-\frac{15}{4}\) & (C) \quad \(\frac{21}{4}\) & (D) \quad \(-\frac{9}{2}\) \\[10pt]
\end{tabularx}}
\item If \(\displaystyle \int x^{3}e^{-2x}\,dx=e^{-2x}g(x)+C\), where \(C\) is a constant of integration and \(g(x)\) is a polynomial, then the value of \(g(0)\) is \hfill [4 marks]
{\renewcommand{\arraystretch}{1.8}
\begin{tabularx}{\linewidth}{X X X X}
(A) \quad \(-\dfrac{3}{8}\) & (B) \quad \(-\dfrac{3}{4}\) & (C) \quad \(\dfrac{3}{8}\) & (D) \quad \(\dfrac{1}{8}\) \\[10pt]
\end{tabularx}}
\item If \(\displaystyle \int x^{4}e^{-2x}\,dx=e^{-2x}g(x)+C\), where \(C\) is a constant of integration and \(g(x)\) is a polynomial, then the value of \(g(1)\) is \hfill [4 marks]
{\renewcommand{\arraystretch}{1.8}
\begin{tabularx}{\linewidth}{X X X X}
(A) \quad \(-\dfrac{21}{4}\) & (B) \quad \(-\dfrac{15}{4}\) & (C) \quad \(\dfrac{3}{4}\) & (D) \quad \(-\dfrac{9}{4}\) \\[10pt]
\end{tabularx}}
\item Let $f(x)$ be a polynomial of degree at most $4$ such that
\[
\int x^4 e^{-2x}\,dx=e^{-2x}f(x)+C,
\]
where $C$ is a constant. Then the value of $f(1)$ is \hfill [4 marks]
{\renewcommand{\arraystretch}{1.8}
\begin{tabularx}{\linewidth}{X X X X}
(A) \quad $-\dfrac{17}{4}$ & (B) \quad $-\dfrac{21}{4}$ & (C) \quad $-\dfrac{23}{4}$ & (D) \quad $-\dfrac{25}{4}$ \\[10pt]
\end{tabularx}}
\item Let \(\displaystyle \int x^{4}e^{-2x}\,dx=e^{-2x}g(x)+C\), where \(C\) is a constant of integration and \(g(x)\) is a polynomial. Then the value of \(g(1)\) is \hfill [4 marks]
{\renewcommand{\arraystretch}{1.8}
\begin{tabularx}{\linewidth}{X X X X}
(A) \quad \(-\dfrac{15}{4}\) & (B) \quad \(-\dfrac{21}{4}\) & (C) \quad \(\dfrac{21}{4}\) & (D) \quad \(-\dfrac{25}{4}\) \\[10pt]
\end{tabularx}}
\item Suppose that
\[
\int x^{4}e^{-2x}\,dx=e^{-2x}g(x)+C,
\]
where \(g(x)\) is a polynomial of degree \(4\) and \(C\) is an arbitrary constant. Then the value of \(g(1)\) is \hfill [4 marks]
{\renewcommand{\arraystretch}{1.8}
\begin{tabularx}{\linewidth}{X X X X}
(A) \quad \(-\dfrac{21}{4}\) & (B) \quad \(-\dfrac{15}{4}\) & (C) \quad \(\dfrac{21}{4}\) & (D) \quad \(-\dfrac{9}{4}\) \\[10pt]
\end{tabularx}}
\item If \(\displaystyle \int x^{4}e^{-2x}\,dx=e^{-2x}g(x)+C\), where \(C\) is a constant of integration and \(g(x)\) is a polynomial, then the value of \(g(1)\) is \hfill [3 marks]
{\renewcommand{\arraystretch}{1.8}
\begin{tabularx}{\linewidth}{X X X X}
(A) \quad \(-\dfrac{21}{4}\) & (B) \quad \(-\dfrac{15}{4}\) & (C) \quad \(\dfrac{21}{4}\) & (D) \quad \(-\dfrac{9}{4}\) \\[10pt]
\end{tabularx}}
\item If \(\displaystyle \int x^{4}e^{-2x}\,dx=e^{-2x}g(x)+C\), where \(g(x)\) is a polynomial and \(C\) is a constant of integration, then the constant term of \(g(x)\) is \hfill [4 marks]
{\renewcommand{\arraystretch}{1.8}
\begin{tabularx}{\linewidth}{X X X X}
(A) \quad \(-\dfrac{3}{4}\) & (B) \quad \(-\dfrac{3}{2}\) & (C) \quad \(\dfrac{3}{4}\) & (D) \quad \(\dfrac{3}{2}\) \\[10pt]
\end{tabularx}}
\item Let \(g(x)\) be a polynomial of degree \(4\) such that
\[
\int x^4e^{-2x}\,dx=e^{-2x}g(x)+C,
\]
where \(C\) is a constant of integration. Then the value of \(g(1)\) is \hfill [4 marks]
{\renewcommand{\arraystretch}{1.8}
\begin{tabularx}{\linewidth}{X X X X}
(A) \quad \(-\dfrac{21}{4}\) & (B) \quad \(-\dfrac{15}{4}\) & (C) \quad \(\dfrac{21}{4}\) & (D) \quad \(\dfrac{15}{4}\) \\[10pt]
\end{tabularx}}
\item Let \[\int x^{4}e^{-2x}\,dx=e^{-2x}g(x)+C,\] where \(C\) is an arbitrary constant and \(g(x)\) is a polynomial. Then the value of \(g(1)\) is \hfill [4 marks]
{\renewcommand{\arraystretch}{1.8}
\begin{tabularx}{\linewidth}{X X X X}
(A) \quad \(-\dfrac{17}{4}\) & (B) \quad \(-\dfrac{21}{4}\) & (C) \quad \(-\dfrac{23}{4}\) & (D) \quad \(-\dfrac{25}{4}\) \\[10pt]
\end{tabularx}}
\item Let \(g(x)\) be a polynomial such that
\[
\int x^4e^{-2x}\,dx=e^{-2x}g(x)+C,
\]
where \(C\) is an arbitrary constant. The value of \(g(1)\) is \hfill [4 marks]
{\renewcommand{\arraystretch}{1.8}
\begin{tabularx}{\linewidth}{X X X X}
(A) \quad \(-\frac{21}{4}\) & (B) \quad \(-\frac{19}{4}\) & (C) \quad \(\frac{21}{4}\) & (D) \quad \(-\frac{25}{4}\) \\[10pt]
\end{tabularx}}
\item Suppose that
\[
\int x^{4}e^{-2x}\,dx=e^{-2x}f(x)+C,
\]
where \(C\) is a constant of integration and \(f(x)\) is a polynomial. Then the value of \(f(1)\) is \hfill [4 marks]
{\renewcommand{\arraystretch}{1.8}
\begin{tabularx}{\linewidth}{X X X X}
(A) \quad \(-\dfrac{21}{4}\) & (B) \quad \(-\dfrac{15}{4}\) & (C) \quad \(\dfrac{21}{4}\) & (D) \quad \(-\dfrac{25}{4}\) \\[10pt]
\end{tabularx}}
\item If \(\displaystyle \int x^{4}e^{-2x}\,dx=e^{-2x}g(x)+C\), where \(C\) is an arbitrary constant and \(g(x)\) is a polynomial, then the value of \(g(0)\) is \hfill [3 marks]
{\renewcommand{\arraystretch}{1.8}
\begin{tabularx}{\linewidth}{X X X X}
(A) \quad \(-\dfrac{3}{4}\) & (B) \quad \(\dfrac{3}{4}\) & (C) \quad \(-\dfrac{1}{2}\) & (D) \quad \(\dfrac{1}{2}\) \\[10pt]
\end{tabularx}}
\item Let $g(x)$ be a polynomial such that
\[
\int x^{4}e^{-2x}\,dx=e^{-2x}g(x)+C,
\]
where $C$ is a constant of integration. The value of $g(-1)$ is \hfill [4 marks]
{\renewcommand{\arraystretch}{1.8}
\begin{tabularx}{\linewidth}{X X X X}
(A) \quad $-\dfrac{1}{4}$ & (B) \quad $\dfrac{1}{4}$ & (C) \quad $-\dfrac{3}{4}$ & (D) \quad $\dfrac{3}{4}$ \\[10pt]
\end{tabularx}}
\item Let \(f(x)\) be a polynomial such that
\[
\int x^4e^{-2x}\,dx=e^{-2x}f(x)+C,
\]
where \(C\) is an arbitrary constant. Then the value of \(f(1)\) is \hfill [4 marks]
{\renewcommand{\arraystretch}{1.8}
\begin{tabularx}{\linewidth}{X X X X}
(A) \quad \(-\frac{21}{4}\) & (B) \quad \(-\frac{19}{4}\) & (C) \quad \(-\frac{23}{4}\) & (D) \quad \(\frac{21}{4}\) \\[10pt]
\end{tabularx}}
\end{enumerate}
\end{document}